% GENERATED FILE. Edit canonical lessons, then run npm run generate:books.
% canonical-source-hash: fnv1a64:05b39b3a3fd11bd5
% canonical-lessons: FA-C458-dastforush, FA-C458-attari, FA-C458-bazzaz, FA-C458-omdeforush, FA-C458-khordeforush

\chapter{Sellers}
\label{ch:fa-sellers}
\hlchaptermodality{458}

\begin{chapteropening}
\textbf{By the end of this chapter:} \emph{I can talk about who sells what: a street vendor, a herb shop, a draper, a wholesaler and a retailer.}

Complete the last lesson of chapter 458: I can talk about who sells what: a street vendor, a herb shop, a draper, a wholesaler and a retailer.
\end{chapteropening}

\section[\texorpdfstring{\fa{دستفروش} (dastforush)}{دستفروش (dastforush)}]{\fa{دستفروش} (dastforush) — a street vendor, a hawker}

\label{lesson:FA-C458-dastforush}



\begin{quote}
Before the new one: say the Persian for a covered hall in the bazaar, then the Persian for a caravanserai.
\end{quote}

\subsection*{You'll want to know: \fa{دستفروش}}
\textbf{\fa{دستفروش}} (\emph{dastforush}) — \textquotedblleft{}a street vendor, a hawker\textquotedblright{}.

\begin{quote}
\fa{در} \fa{مترو} \fa{یک} \fa{دستفروش} \fa{جوراب} \fa{میفروخت}.
\end{quote}

\begin{quote}
\emph{Dar metro yek dastforush jurâb miforukht.}
\end{quote}

\begin{quote}
A hawker was selling socks in the metro.
\end{quote}

Literally 'hand-seller': \emph{dast} 'hand' + \emph{forush} 'seller'.

\begin{cousinweb}
\emph{dastforush} is a native compound of \emph{dast} 'hand', which EtymDB traces to Old Persian \emph{dasta-}, and \emph{forush}, the present stem of \emph{forukhtan} 'to sell', usually traced to Middle Persian \emph{frōxtan}.
\end{cousinweb}

\begin{grammarlens}[title={What you've built}]
One more word for telling, explaining and describing at B1.
\end{grammarlens}

\subsection*{Your turn}
\begin{itemize}
\raggedright
  \item \emph{Say it:} \emph{dastforush}
  \item \emph{Say it:} \emph{dastforush}, once more
  \item \emph{Say it:} a sentence of your own with \emph{dastforush}
  \item \emph{Recall:} say the Persian for a covered hall in the bazaar, then the Persian for a caravanserai, then say \emph{dastforush} again
\end{itemize}

\begin{tcolorbox}[breakable,colback=teal!4,colframe=teal!35!black,title={Before you move on}]
What does \fa{دستفروش} mean? (\textquotedblleft{}A street vendor, a hawker\textquotedblright{}.) Say it once more.
\end{tcolorbox}
\section[\texorpdfstring{\fa{عطاری} ('attâri)}{عطاری ('attâri)}]{\fa{عطاری} ('attâri) — a herbalist's shop, a herb and spice shop}

\label{lesson:FA-C458-attari}



\begin{quote}
Before the new one: say the Persian for a caravanserai, then the Persian for a street vendor.
\end{quote}

\subsection*{You'll want to know: \fa{عطاری}}
\textbf{\fa{عطاری}} (\emph{'attâri}) — \textquotedblleft{}a herbalist's shop, a herb and spice shop\textquotedblright{}.

\begin{quote}
\fa{مادرم} \fa{از} \fa{عطاری} \fa{دارچین} \fa{خرید}.
\end{quote}

\begin{quote}
\emph{Mâdaram az 'attâri dârchin kharid.}
\end{quote}

\begin{quote}
My mother bought cinnamon at the herbalist's.
\end{quote}

The seller is the \emph{'attâr}. These shops sell dried herbs, spices and old remedies.

\begin{cousinweb}
\emph{'attâri} is built from \emph{'attâr} 'a seller of perfumes and herbs', usually traced to Arabic \emph{ʿaṭṭār}, a word built on \emph{ʿiṭr} 'perfume', which EtymDB lists, with the Persian ending \emph{-i}.
\end{cousinweb}

\begin{grammarlens}[title={What you've built}]
One more word for telling, explaining and describing at B1.
\end{grammarlens}

\subsection*{Your turn}
\begin{itemize}
\raggedright
  \item \emph{Say it:} \emph{'attâri}
  \item \emph{Say it:} \emph{'attâri}, once more
  \item \emph{Say it:} a sentence of your own with \emph{'attâri}
  \item \emph{Recall:} say the Persian for a caravanserai, then the Persian for a street vendor, then say \emph{'attâri} again
\end{itemize}

\begin{tcolorbox}[breakable,colback=teal!4,colframe=teal!35!black,title={Before you move on}]
What does \fa{عطاری} mean? (\textquotedblleft{}A herbalist's shop, a herb and spice shop\textquotedblright{}.) Say it once more.
\end{tcolorbox}
\section[\texorpdfstring{\fa{بزاز} (bazzâz)}{بزاز (bazzâz)}]{\fa{بزاز} (bazzâz) — a draper, a cloth merchant}

\label{lesson:FA-C458-bazzaz}



\begin{quote}
Before the new one: say the Persian for a street vendor, then the Persian for a herbalist's shop.
\end{quote}

\subsection*{You'll want to know: \fa{بزاز}}
\textbf{\fa{بزاز}} (\emph{bazzâz}) — \textquotedblleft{}a draper, a cloth merchant\textquotedblright{}.

\begin{quote}
\fa{بزاز} \fa{پارچه} \fa{را} \fa{با} \fa{متر} \fa{اندازه} \fa{گرفت}.
\end{quote}

\begin{quote}
\emph{Bazzâz pârche râ bâ metr andâze gereft.}
\end{quote}

\begin{quote}
The draper measured the cloth with a tape.
\end{quote}

An old word; today often \emph{pârche-forush}.

\begin{cousinweb}
\emph{bazzâz} was borrowed from Arabic \emph{bazzāz} 'draper, cloth merchant' (EtymDB).
\end{cousinweb}

\begin{grammarlens}[title={What you've built}]
One more word for telling, explaining and describing at B1.
\end{grammarlens}

\subsection*{Your turn}
\begin{itemize}
\raggedright
  \item \emph{Say it:} \emph{bazzâz}
  \item \emph{Say it:} \emph{bazzâz}, once more
  \item \emph{Say it:} a sentence of your own with \emph{bazzâz}
  \item \emph{Recall:} say the Persian for a street vendor, then the Persian for a herbalist's shop, then say \emph{bazzâz} again
\end{itemize}

\begin{tcolorbox}[breakable,colback=teal!4,colframe=teal!35!black,title={Before you move on}]
What does \fa{بزاز} mean? (\textquotedblleft{}A draper, a cloth merchant\textquotedblright{}.) Say it once more.
\end{tcolorbox}
\section[\texorpdfstring{\fa{عمده} \fa{فروش} ('omdeforush)}{عمده فروش ('omdeforush)}]{\fa{عمده} \fa{فروش} ('omdeforush) — a wholesaler}

\label{lesson:FA-C458-omdeforush}



\begin{quote}
Before the new one: say the Persian for a herbalist's shop, then the Persian for a draper.
\end{quote}

\subsection*{You'll want to know: \fa{عمده} \fa{فروش}}
\textbf{\fa{عمده} \fa{فروش}} (\emph{'omdeforush}) — \textquotedblleft{}a wholesaler\textquotedblright{}.

\begin{quote}
\fa{عمده} \fa{فروش} \fa{برنج} \fa{را} \fa{کیسه} \fa{ای} \fa{میفروشد}.
\end{quote}

\begin{quote}
\emph{'Omdeforush berenj râ kise-i miforushad.}
\end{quote}

\begin{quote}
The wholesaler sells rice by the sack.
\end{quote}

\emph{'omde} 'in bulk' + \emph{forush} 'seller'.

\begin{cousinweb}
\emph{'omdeforush} is built from \emph{'omde} 'main part, in bulk', which EtymDB takes from Arabic \emph{ʿumda}, and native \emph{forush}, the present stem of \emph{forukhtan} 'to sell'.
\end{cousinweb}

\begin{grammarlens}[title={What you've built}]
One more word for telling, explaining and describing at B1.
\end{grammarlens}

\subsection*{Your turn}
\begin{itemize}
\raggedright
  \item \emph{Say it:} \emph{'omdeforush}
  \item \emph{Say it:} \emph{'omdeforush}, once more
  \item \emph{Say it:} a sentence of your own with \emph{'omdeforush}
  \item \emph{Recall:} say the Persian for a herbalist's shop, then the Persian for a draper, then say \emph{'omdeforush} again
\end{itemize}

\begin{tcolorbox}[breakable,colback=teal!4,colframe=teal!35!black,title={Before you move on}]
What does \fa{عمده} \fa{فروش} mean? (\textquotedblleft{}A wholesaler\textquotedblright{}.) Say it once more.
\end{tcolorbox}
\section[\texorpdfstring{\fa{خرده} \fa{فروش} (khordeforush)}{خرده فروش (khordeforush)}]{\fa{خرده} \fa{فروش} (khordeforush) — a retailer}

\label{lesson:FA-C458-khordeforush}



\begin{quote}
Before the new one: say the Persian for a draper, then the Persian for a wholesaler.
\end{quote}

\subsection*{You'll want to know: \fa{خرده} \fa{فروش}}
\textbf{\fa{خرده} \fa{فروش}} (\emph{khordeforush}) — \textquotedblleft{}a retailer\textquotedblright{}.

\begin{quote}
\fa{خرده} \fa{فروش} \fa{برنج} \fa{را} \fa{کیلویی} \fa{میفروشد}.
\end{quote}

\begin{quote}
\emph{Khordeforush berenj râ kiluyi miforushad.}
\end{quote}

\begin{quote}
The retailer sells rice by the kilo.
\end{quote}

\emph{khorde} 'small bits' + \emph{forush} 'seller'; the opposite of \emph{'omdeforush}.

\begin{cousinweb}
\emph{khordeforush} is a native compound of \emph{khorde} 'small bits', which EtymDB lists, and \emph{forush}, the present stem of \emph{forukhtan} 'to sell'.
\end{cousinweb}

\begin{grammarlens}[title={What you've built}]
One more word for telling, explaining and describing at B1.
\end{grammarlens}

\subsection*{Your turn}
\begin{itemize}
\raggedright
  \item \emph{Say it:} \emph{khordeforush}
  \item \emph{Say it:} \emph{khordeforush}, once more
  \item \emph{Say it:} a sentence of your own with \emph{khordeforush}
  \item \emph{Recall:} say the Persian for a draper, then the Persian for a wholesaler, then say \emph{khordeforush} again
\end{itemize}

\begin{tcolorbox}[breakable,colback=teal!4,colframe=teal!35!black,title={Before you move on}]
What does \fa{خرده} \fa{فروش} mean? (\textquotedblleft{}A retailer\textquotedblright{}.) Say it once more.
\end{tcolorbox}
